% Options for packages loaded elsewhere
\PassOptionsToPackage{unicode}{hyperref}
\PassOptionsToPackage{hyphens}{url}
%
\documentclass[
  11pt,
]{article}
\usepackage{amsmath,amssymb}
\usepackage{iftex}
\ifPDFTeX
  \usepackage[T1]{fontenc}
  \usepackage[utf8]{inputenc}
  \usepackage{textcomp} % provide euro and other symbols
\else % if luatex or xetex
  \usepackage{unicode-math} % this also loads fontspec
  \defaultfontfeatures{Scale=MatchLowercase}
  \defaultfontfeatures[\rmfamily]{Ligatures=TeX,Scale=1}
\fi
\usepackage{lmodern}
\ifPDFTeX\else
  % xetex/luatex font selection
\fi
% Use upquote if available, for straight quotes in verbatim environments
\IfFileExists{upquote.sty}{\usepackage{upquote}}{}
\IfFileExists{microtype.sty}{% use microtype if available
  \usepackage[]{microtype}
  \UseMicrotypeSet[protrusion]{basicmath} % disable protrusion for tt fonts
}{}
\makeatletter
\@ifundefined{KOMAClassName}{% if non-KOMA class
  \IfFileExists{parskip.sty}{%
    \usepackage{parskip}
  }{% else
    \setlength{\parindent}{0pt}
    \setlength{\parskip}{6pt plus 2pt minus 1pt}}
}{% if KOMA class
  \KOMAoptions{parskip=half}}
\makeatother
\usepackage{xcolor}
\usepackage[margin=.5in]{geometry}
\usepackage{color}
\usepackage{fancyvrb}
\newcommand{\VerbBar}{|}
\newcommand{\VERB}{\Verb[commandchars=\\\{\}]}
\DefineVerbatimEnvironment{Highlighting}{Verbatim}{commandchars=\\\{\}}
% Add ',fontsize=\small' for more characters per line
\usepackage{framed}
\definecolor{shadecolor}{RGB}{248,248,248}
\newenvironment{Shaded}{\begin{snugshade}}{\end{snugshade}}
\newcommand{\AlertTok}[1]{\textcolor[rgb]{0.94,0.16,0.16}{#1}}
\newcommand{\AnnotationTok}[1]{\textcolor[rgb]{0.56,0.35,0.01}{\textbf{\textit{#1}}}}
\newcommand{\AttributeTok}[1]{\textcolor[rgb]{0.13,0.29,0.53}{#1}}
\newcommand{\BaseNTok}[1]{\textcolor[rgb]{0.00,0.00,0.81}{#1}}
\newcommand{\BuiltInTok}[1]{#1}
\newcommand{\CharTok}[1]{\textcolor[rgb]{0.31,0.60,0.02}{#1}}
\newcommand{\CommentTok}[1]{\textcolor[rgb]{0.56,0.35,0.01}{\textit{#1}}}
\newcommand{\CommentVarTok}[1]{\textcolor[rgb]{0.56,0.35,0.01}{\textbf{\textit{#1}}}}
\newcommand{\ConstantTok}[1]{\textcolor[rgb]{0.56,0.35,0.01}{#1}}
\newcommand{\ControlFlowTok}[1]{\textcolor[rgb]{0.13,0.29,0.53}{\textbf{#1}}}
\newcommand{\DataTypeTok}[1]{\textcolor[rgb]{0.13,0.29,0.53}{#1}}
\newcommand{\DecValTok}[1]{\textcolor[rgb]{0.00,0.00,0.81}{#1}}
\newcommand{\DocumentationTok}[1]{\textcolor[rgb]{0.56,0.35,0.01}{\textbf{\textit{#1}}}}
\newcommand{\ErrorTok}[1]{\textcolor[rgb]{0.64,0.00,0.00}{\textbf{#1}}}
\newcommand{\ExtensionTok}[1]{#1}
\newcommand{\FloatTok}[1]{\textcolor[rgb]{0.00,0.00,0.81}{#1}}
\newcommand{\FunctionTok}[1]{\textcolor[rgb]{0.13,0.29,0.53}{\textbf{#1}}}
\newcommand{\ImportTok}[1]{#1}
\newcommand{\InformationTok}[1]{\textcolor[rgb]{0.56,0.35,0.01}{\textbf{\textit{#1}}}}
\newcommand{\KeywordTok}[1]{\textcolor[rgb]{0.13,0.29,0.53}{\textbf{#1}}}
\newcommand{\NormalTok}[1]{#1}
\newcommand{\OperatorTok}[1]{\textcolor[rgb]{0.81,0.36,0.00}{\textbf{#1}}}
\newcommand{\OtherTok}[1]{\textcolor[rgb]{0.56,0.35,0.01}{#1}}
\newcommand{\PreprocessorTok}[1]{\textcolor[rgb]{0.56,0.35,0.01}{\textit{#1}}}
\newcommand{\RegionMarkerTok}[1]{#1}
\newcommand{\SpecialCharTok}[1]{\textcolor[rgb]{0.81,0.36,0.00}{\textbf{#1}}}
\newcommand{\SpecialStringTok}[1]{\textcolor[rgb]{0.31,0.60,0.02}{#1}}
\newcommand{\StringTok}[1]{\textcolor[rgb]{0.31,0.60,0.02}{#1}}
\newcommand{\VariableTok}[1]{\textcolor[rgb]{0.00,0.00,0.00}{#1}}
\newcommand{\VerbatimStringTok}[1]{\textcolor[rgb]{0.31,0.60,0.02}{#1}}
\newcommand{\WarningTok}[1]{\textcolor[rgb]{0.56,0.35,0.01}{\textbf{\textit{#1}}}}
\usepackage{longtable,booktabs,array}
\usepackage{calc} % for calculating minipage widths
% Correct order of tables after \paragraph or \subparagraph
\usepackage{etoolbox}
\makeatletter
\patchcmd\longtable{\par}{\if@noskipsec\mbox{}\fi\par}{}{}
\makeatother
% Allow footnotes in longtable head/foot
\IfFileExists{footnotehyper.sty}{\usepackage{footnotehyper}}{\usepackage{footnote}}
\makesavenoteenv{longtable}
\usepackage{graphicx}
\makeatletter
\def\maxwidth{\ifdim\Gin@nat@width>\linewidth\linewidth\else\Gin@nat@width\fi}
\def\maxheight{\ifdim\Gin@nat@height>\textheight\textheight\else\Gin@nat@height\fi}
\makeatother
% Scale images if necessary, so that they will not overflow the page
% margins by default, and it is still possible to overwrite the defaults
% using explicit options in \includegraphics[width, height, ...]{}
\setkeys{Gin}{width=\maxwidth,height=\maxheight,keepaspectratio}
% Set default figure placement to htbp
\makeatletter
\def\fps@figure{htbp}
\makeatother
\setlength{\emergencystretch}{3em} % prevent overfull lines
\providecommand{\tightlist}{%
  \setlength{\itemsep}{0pt}\setlength{\parskip}{0pt}}
\setcounter{secnumdepth}{-\maxdimen} % remove section numbering
\usepackage{amsmath}
\usepackage{mathtools}
\usepackage{amsthm}
\usepackage{bm}
\usepackage{xcolor}
\usepackage{fbox}
\usepackage{amsmath}
\usepackage{hw4}
\ifLuaTeX
  \usepackage{selnolig}  % disable illegal ligatures
\fi
\IfFileExists{bookmark.sty}{\usepackage{bookmark}}{\usepackage{hyperref}}
\IfFileExists{xurl.sty}{\usepackage{xurl}}{} % add URL line breaks if available
\urlstyle{same}
\hypersetup{
  pdftitle={STAT 581 - Problem Set 6},
  pdfauthor={Alex Towell (atowell@siue.edu)},
  hidelinks,
  pdfcreator={LaTeX via pandoc}}

\title{STAT 581 - Problem Set 6}
\author{Alex Towell
(\href{mailto:atowell@siue.edu}{\nolinkurl{atowell@siue.edu}})}
\date{}

\begin{document}
\maketitle

\hypertarget{problem-1}{%
\section{Problem 1}\label{problem-1}}

A soft drink bottler is interested in studying the effects on a filling
process. A factorial experiment is run using three factors: percent
carbonation (in \%), operating pressure (in psi), and line speed (in
bpm). The response variable is deviation from the target fill height.
The experiment consists of \(n=2\) runs for each of the \(a b c = 12\)
treatment combinations. The data is available on Blackboard as an Excel
File.

\hypertarget{preliminary-analysis}{%
\subsection{Preliminary analysis}\label{preliminary-analysis}}

The three-factor factorial experiment with \(a=3\) levels of the
\emph{carbonation} factor (in \%), denoted by \(A\), \(b=2\) levels of
the operating \emph{pressure} (in psi), denoted by \(B\), and \(c=2\)
levels of the line \emph{speed} (in bpm), denoted by \(C\). Each
treatment combination has \(n=2\) replicates for a total of
\(N=a b c n=24\) observations. The response variable \(y\) is deviation
from the target fill height.

We describe the observations in the three-factor factorial experiment by
the fixed-effects interaction model where treatments are defined to be
deviations from the overall mean, \[
  Y_{i j k \ell} = \mu + \tau_i + \beta_j + \gamma_k +
                   (\tau\!\beta)_{i j} + (\tau\!\gamma)_{i k} +
                   (\beta\!\gamma)_{j k} + (\tau\!\beta\gamma)_{i j k} +
                   \epsilon_{i j k \ell}
  \begin{cases}
    i = 1,\ldots,a=3\\
    j = 1,\ldots,b=2\\
    k = 1,\ldots,c=2\\
    \ell = 1,\ldots,n=2\\
  \end{cases}
\] where \(\mu\) is the overall main effect, \(\tau_i\) is the effect of
the \(i\)-th level of factor \(A\), \(\beta_j\) is the effect of the
\(j\)-th level of factor \(B\), and \(\gamma_k\) is the effect of the
\(k\)-th level of factor \(C\). The interaction effects may also be
observed using the same labeling, e.g., \((\tau\!\gamma)_{i k}\) is the
interaction effect of the \(i\)-th level of factor \(A\) and the
\(k\)-th level of factor \(C\).

\hypertarget{a}{%
\subsection{(a)}\label{a}}

\fbox{\begin{minipage}{\textwidth}
Perform tests for all main effects and for all interaction effects.
State the $F$-statistic and $p$-value for each of the main effect tests, and for
the test of an interaction between carbonation and pressure.
\end{minipage}}

\begin{Shaded}
\begin{Highlighting}[]
\FunctionTok{library}\NormalTok{(}\StringTok{"readxl"}\NormalTok{)}
\NormalTok{h6.data }\OtherTok{=} \FunctionTok{read\_excel}\NormalTok{(}\StringTok{"handout6data.xlsx"}\NormalTok{)}

\NormalTok{A }\OtherTok{=} \FunctionTok{as.factor}\NormalTok{(}\FunctionTok{na.omit}\NormalTok{(h6.data}\SpecialCharTok{$}\NormalTok{carbon))}
\NormalTok{B }\OtherTok{=} \FunctionTok{as.factor}\NormalTok{(}\FunctionTok{na.omit}\NormalTok{(h6.data}\SpecialCharTok{$}\NormalTok{pressure))}
\NormalTok{C }\OtherTok{=} \FunctionTok{as.factor}\NormalTok{(}\FunctionTok{na.omit}\NormalTok{(h6.data}\SpecialCharTok{$}\NormalTok{speed))}
\NormalTok{y }\OtherTok{=} \FunctionTok{na.omit}\NormalTok{(h6.data}\SpecialCharTok{$}\NormalTok{deviation)}

\FunctionTok{contrasts}\NormalTok{(A) }\OtherTok{=}\NormalTok{ contr.sum }\CommentTok{\# we use these in the additional analysis section}
\FunctionTok{contrasts}\NormalTok{(B) }\OtherTok{=}\NormalTok{ contr.sum }\CommentTok{\# to make it easier to provide explicit estimates of}
\FunctionTok{contrasts}\NormalTok{(C) }\OtherTok{=}\NormalTok{ contr.sum }\CommentTok{\# the parameters}

\NormalTok{x.mod }\OtherTok{=} \FunctionTok{aov}\NormalTok{(y }\SpecialCharTok{\textasciitilde{}}\NormalTok{ A}\SpecialCharTok{*}\NormalTok{B}\SpecialCharTok{*}\NormalTok{C)}
\FunctionTok{summary}\NormalTok{(x.mod)}
\end{Highlighting}
\end{Shaded}

\begin{verbatim}
##             Df Sum Sq Mean Sq F value   Pr(>F)    
## A            2 252.75  126.38 178.412 1.19e-09 ***
## B            1  45.37   45.37  64.059 3.74e-06 ***
## C            1  22.04   22.04  31.118  0.00012 ***
## A:B          2   5.25    2.63   3.706  0.05581 .  
## A:C          2   0.58    0.29   0.412  0.67149    
## B:C          1   1.04    1.04   1.471  0.24859    
## A:B:C        2   1.08    0.54   0.765  0.48687    
## Residuals   12   8.50    0.71                     
## ---
## Signif. codes:  0 '***' 0.001 '**' 0.01 '*' 0.05 '.' 0.1 ' ' 1
\end{verbatim}

The tests for each of the main effects and a test of an interaction
between carbon \((A\)) and pressure \((B\)) is given by the following
table:

\begin{longtable}[]{@{}lll@{}}
\toprule\noalign{}
effects & \(F^*\) & \(p\)-value \\
\midrule\noalign{}
\endhead
\bottomrule\noalign{}
\endlastfoot
\(A\) & \(178.4\) & \(.000\) \\
\(B\) & \(64.1\) & \(.000\) \\
\(C\) & \(31.1\) & \(.000\) \\
\(A \times B\) & \(3.7\) & \(.056\) \\
\end{longtable}

\hypertarget{b}{%
\subsection{(b)}\label{b}}

\fbox{\begin{minipage}{\textwidth}
Explain how an interaction plot is used in studying an interaction effect.
\end{minipage}}

If an interaction plot is parallel, or nearly so, then there is no need
to include the interaction term.

In the \(3\)-factor ANOVA model, a \(3\)-factor interaction occurs when
a \(2\)-factor interaction depends on the level of the \(3\)rd factor.

\hypertarget{c}{%
\subsection{(c)}\label{c}}

\fbox{\begin{minipage}{\textwidth}
Determine fitted values for the main effects model, as well as fitted values for
the model including the carbonation / pressure interaction.
Plot the fitted values for each model using an interaction plot.
Comment on the decision of whether or not to include the interaction term.
\end{minipage}}

From the earlier results, we saw that factors \(A\) (carbon), \(B\)
(pressure), and \(C\) (speed) were statistically significant, but only
the interaction between factor \(A\) and \(B\) (carbon and pressure) had
any significance.

\hypertarget{additive-model}{%
\subsubsection{Additive model}\label{additive-model}}

The additive model is given by \[
  Y_{i j k \ell}^{(R)} = \mu + \tau_i + \beta_j + \gamma_k + \epsilon_{i j k \ell}.
\] Thus, \[
  E(Y_{i j k}^{(R)}) = \mu + \tau_i + \beta_j + \gamma_k.
\]

We fit the additive model with:

\begin{Shaded}
\begin{Highlighting}[]
\NormalTok{mod.R }\OtherTok{=} \FunctionTok{aov}\NormalTok{(y }\SpecialCharTok{\textasciitilde{}}\NormalTok{ A}\SpecialCharTok{+}\NormalTok{B}\SpecialCharTok{+}\NormalTok{C)}
\end{Highlighting}
\end{Shaded}

We define the estimated means for the models with:

\begin{Shaded}
\begin{Highlighting}[]
\NormalTok{fits.R }\OtherTok{=} \FunctionTok{predict}\NormalTok{(mod.R)}
\end{Highlighting}
\end{Shaded}

The interaction plots for the additive where the left and right plots
are respectively for factor \(C\) at levels \(1\) and \(2\).

\begin{Shaded}
\begin{Highlighting}[]
\FunctionTok{par}\NormalTok{(}\AttributeTok{mfrow=}\FunctionTok{c}\NormalTok{(}\DecValTok{1}\NormalTok{,}\DecValTok{2}\NormalTok{))}
\ControlFlowTok{for}\NormalTok{ (i }\ControlFlowTok{in} \DecValTok{1}\SpecialCharTok{:}\FunctionTok{nlevels}\NormalTok{(C))}
\NormalTok{\{}
  \FunctionTok{interaction.plot}\NormalTok{(A[C}\SpecialCharTok{==}\FunctionTok{levels}\NormalTok{(C)[i]],B[C}\SpecialCharTok{==}\FunctionTok{levels}\NormalTok{(C)[i]],}
\NormalTok{                   fits.R[C}\SpecialCharTok{==}\FunctionTok{levels}\NormalTok{(C)[i]],}
                   \AttributeTok{xlab=}\StringTok{"A"}\NormalTok{,}
                   \AttributeTok{ylab=}\StringTok{"mean response"}\NormalTok{,}
                   \AttributeTok{trace.label=}\StringTok{"B"}\NormalTok{)}
\NormalTok{\}}
\end{Highlighting}
\end{Shaded}

\includegraphics{hw6_files/figure-latex/unnamed-chunk-5-1.pdf}

\hypertarget{interaction-model}{%
\subsubsection{Interaction model}\label{interaction-model}}

The model that includes interactions between factors \(A\) (carbon) and
\(B\) (pressure) is given by \[
  Y_{i j k \ell}^{(F)} = \mu + \tau_i + \beta_j + \gamma_k +
                   (\tau\!\beta)_{i j} + \epsilon_{i j k \ell}.
\] Thus, \[
  E(Y_{i j k}^{(F)}) = \mu + \tau_i + \beta_j + \gamma_k + (\tau\!\beta)_{i j}.
\]

We fit the interaction model with:

\begin{Shaded}
\begin{Highlighting}[]
\NormalTok{mod.F }\OtherTok{=} \FunctionTok{aov}\NormalTok{(y }\SpecialCharTok{\textasciitilde{}}\NormalTok{ A}\SpecialCharTok{+}\NormalTok{B}\SpecialCharTok{+}\NormalTok{C}\SpecialCharTok{+}\NormalTok{A}\SpecialCharTok{:}\NormalTok{B)}
\end{Highlighting}
\end{Shaded}

We define the estimated means for the models with:

\begin{Shaded}
\begin{Highlighting}[]
\NormalTok{fits.F }\OtherTok{=} \FunctionTok{predict}\NormalTok{(mod.F)}
\end{Highlighting}
\end{Shaded}

The interaction plots for the interaction model where the left and right
plots are respectively for factor \(C\) at levels \(1\) and \(2\).

\begin{Shaded}
\begin{Highlighting}[]
\FunctionTok{par}\NormalTok{(}\AttributeTok{mfrow=}\FunctionTok{c}\NormalTok{(}\DecValTok{1}\NormalTok{,}\DecValTok{2}\NormalTok{))}
\ControlFlowTok{for}\NormalTok{ (i }\ControlFlowTok{in} \DecValTok{1}\SpecialCharTok{:}\FunctionTok{nlevels}\NormalTok{(C))}
\NormalTok{\{}
  \FunctionTok{interaction.plot}\NormalTok{(A[C}\SpecialCharTok{==}\FunctionTok{levels}\NormalTok{(C)[i]],B[C}\SpecialCharTok{==}\FunctionTok{levels}\NormalTok{(C)[i]],}
\NormalTok{                   fits.F[C}\SpecialCharTok{==}\FunctionTok{levels}\NormalTok{(C)[i]],}
                   \AttributeTok{xlab=}\StringTok{"A"}\NormalTok{,}
                   \AttributeTok{ylab=}\StringTok{"mean response"}\NormalTok{,}
                   \AttributeTok{trace.label=}\StringTok{"B"}\NormalTok{)}
\NormalTok{\}}
\end{Highlighting}
\end{Shaded}

\includegraphics{hw6_files/figure-latex/unnamed-chunk-8-1.pdf}

\hypertarget{decision}{%
\subsubsection{Decision}\label{decision}}

The interaction model fits show only a slight difference difference from
the additive model. As a general rule, we prefer simpler models to
complex models, and so we decide to leave out the interaction term.

\hypertarget{d}{%
\subsection{(d)}\label{d}}

\fbox{\begin{minipage}{\textwidth}
Create box plots showing the main effects.
Provide an overall conclusion, stated in the context of the problem.
\end{minipage}}

\begin{Shaded}
\begin{Highlighting}[]
\FunctionTok{par}\NormalTok{(}\AttributeTok{mfrow=}\FunctionTok{c}\NormalTok{(}\DecValTok{1}\NormalTok{,}\DecValTok{3}\NormalTok{))}
\FunctionTok{plot}\NormalTok{(fits.R }\SpecialCharTok{\textasciitilde{}}\NormalTok{ A)}
\FunctionTok{plot}\NormalTok{(fits.R }\SpecialCharTok{\textasciitilde{}}\NormalTok{ B)}
\FunctionTok{plot}\NormalTok{(fits.R }\SpecialCharTok{\textasciitilde{}}\NormalTok{ C)}
\end{Highlighting}
\end{Shaded}

\includegraphics{hw6_files/figure-latex/unnamed-chunk-9-1.pdf}

We see that each factor has a positive effect on fill height. The factor
with the biggest marginal effect on fill height is the percent of
carbonation (factor \(A\)).

\hypertarget{problem-2}{%
\section{Problem 2}\label{problem-2}}

An engineer is interested in the effects of cutting speed, tool
geometry, and cutting angle on the life (in hours) of a machine tool.
Two levels of each factor are set, and \(n=3\) runs of a \(2^3\) design
are completed. The data is available on Blackboard as an Excel File.

\hypertarget{a-1}{%
\subsection{(a)}\label{a-1}}

\fbox{\begin{minipage}{\textwidth}
Perform tests for all main effects and for all interaction effects.
State the $F$-statistic and $p$-value for each test of an effect deemed to be
important.
\end{minipage}}

\begin{Shaded}
\begin{Highlighting}[]
\NormalTok{data2 }\OtherTok{=} \FunctionTok{read\_excel}\NormalTok{(}\StringTok{"handout6data.xlsx"}\NormalTok{)}

\NormalTok{A }\OtherTok{=} \FunctionTok{as.factor}\NormalTok{(}\FunctionTok{na.omit}\NormalTok{(data2}\SpecialCharTok{$}\NormalTok{spd))}
\NormalTok{B }\OtherTok{=} \FunctionTok{as.factor}\NormalTok{(}\FunctionTok{na.omit}\NormalTok{(data2}\SpecialCharTok{$}\NormalTok{geo))}
\NormalTok{C }\OtherTok{=} \FunctionTok{as.factor}\NormalTok{(}\FunctionTok{na.omit}\NormalTok{(data2}\SpecialCharTok{$}\NormalTok{angle))}
\NormalTok{y }\OtherTok{=} \FunctionTok{na.omit}\NormalTok{(h6.data}\SpecialCharTok{$}\NormalTok{life)}

\FunctionTok{contrasts}\NormalTok{(A) }\OtherTok{=}\NormalTok{ contr.sum}
\FunctionTok{contrasts}\NormalTok{(B) }\OtherTok{=}\NormalTok{ contr.sum}
\FunctionTok{contrasts}\NormalTok{(C) }\OtherTok{=}\NormalTok{ contr.sum}

\NormalTok{saturated.mod }\OtherTok{=} \FunctionTok{aov}\NormalTok{(y }\SpecialCharTok{\textasciitilde{}}\NormalTok{ A}\SpecialCharTok{*}\NormalTok{B}\SpecialCharTok{*}\NormalTok{C)}
\FunctionTok{summary}\NormalTok{(saturated.mod)}
\end{Highlighting}
\end{Shaded}

\begin{verbatim}
##             Df Sum Sq Mean Sq F value   Pr(>F)    
## A            1    0.7     0.7   0.022 0.883680    
## B            1  770.7   770.7  25.547 0.000117 ***
## C            1  280.2   280.2   9.287 0.007679 ** 
## A:B          1   16.7    16.7   0.552 0.468078    
## A:C          1  468.2   468.2  15.519 0.001172 ** 
## B:C          1   48.2    48.2   1.597 0.224475    
## A:B:C        1   28.2    28.2   0.934 0.348282    
## Residuals   16  482.7    30.2                     
## ---
## Signif. codes:  0 '***' 0.001 '**' 0.01 '*' 0.05 '.' 0.1 ' ' 1
\end{verbatim}

Factors \(B\) and \(C\) are important effects and the the interaction
between \(A\) and \(C\) is also important. Thus, \(A\), \(B\), \(C\),
and \(A:C\) should be included in the model (we include \(A\) because of
the importance of \(A:C\)).

\begin{longtable}[]{@{}lll@{}}
\toprule\noalign{}
effects & \(F^*\) & \(p\)-value \\
\midrule\noalign{}
\endhead
\bottomrule\noalign{}
\endlastfoot
\(A\) & \(.022\) & \(.884\) \\
\(B\) & \(25.547\) & \(.000\) \\
\(C\) & \(9.287\) & \(.008\) \\
\(A:C\) & \(15.519\) & \(.001\) \\
\end{longtable}

\hypertarget{b-1}{%
\subsection{(b)}\label{b-1}}

\fbox{\begin{minipage}{\textwidth}
Fit a reduced model with the main effects, and the statistically significant
interaction effect.
\end{minipage}}

\begin{Shaded}
\begin{Highlighting}[]
\NormalTok{reduced.mod }\OtherTok{=} \FunctionTok{aov}\NormalTok{(y }\SpecialCharTok{\textasciitilde{}}\NormalTok{ A}\SpecialCharTok{+}\NormalTok{B}\SpecialCharTok{+}\NormalTok{C}\SpecialCharTok{+}\NormalTok{A}\SpecialCharTok{:}\NormalTok{C)}
\FunctionTok{summary}\NormalTok{(reduced.mod)}
\end{Highlighting}
\end{Shaded}

\begin{verbatim}
##             Df Sum Sq Mean Sq F value   Pr(>F)    
## A            1    0.7     0.7   0.022 0.883641    
## B            1  770.7   770.7  25.436 7.22e-05 ***
## C            1  280.2   280.2   9.247 0.006724 ** 
## A:C          1  468.2   468.2  15.452 0.000897 ***
## Residuals   19  575.7    30.3                     
## ---
## Signif. codes:  0 '***' 0.001 '**' 0.01 '*' 0.05 '.' 0.1 ' ' 1
\end{verbatim}

\hypertarget{i}{%
\subsubsection{(i)}\label{i}}

\fbox{\begin{minipage}{\textwidth}
Create a plot for the interaction effect.
Provide an interpretation, stated in the context of the problem.
\end{minipage}}

\begin{Shaded}
\begin{Highlighting}[]
\CommentTok{\# define the estimated means for the reduced model for plotting (randomness has}
\CommentTok{\# been smoothed)}
\NormalTok{fits.r }\OtherTok{=} \FunctionTok{predict}\NormalTok{(reduced.mod)}
\FunctionTok{interaction.plot}\NormalTok{(A,C,fits.r,}\AttributeTok{ylab=}\StringTok{"y"}\NormalTok{)}
\end{Highlighting}
\end{Shaded}

\includegraphics{hw6_files/figure-latex/unnamed-chunk-12-1.pdf}

\hypertarget{interpretation}{%
\paragraph{Interpretation}\label{interpretation}}

The experiment finds that when cutting speed (factor \(A\)) is low,
cutting angle has a positive effect on the life of a machine tool
(response).

\hypertarget{ii}{%
\subsubsection{(ii)}\label{ii}}

\fbox{\begin{minipage}{\textwidth}
Create a box plot showing the main effect for the remaining factor.
Provide an interpretation, stated in the context of the problem.
\end{minipage}}

\begin{Shaded}
\begin{Highlighting}[]
\FunctionTok{plot}\NormalTok{(fits.r }\SpecialCharTok{\textasciitilde{}}\NormalTok{ B)}
\end{Highlighting}
\end{Shaded}

\includegraphics{hw6_files/figure-latex/unnamed-chunk-13-1.pdf}

\hypertarget{interpetation}{%
\paragraph{Interpetation}\label{interpetation}}

Tool geometry (factor \(B\)) has a positive effect on the life of a
machine tool.

\hypertarget{iii}{%
\subsubsection{(iii)}\label{iii}}

\fbox{\begin{minipage}{\textwidth}
Create a plot of the fitted values for the reduced model.
Which setting of the factors should be used if the goal is to maximize tool
life.
Explain how this result follows from the interaction plot and the main effects
plot.
\end{minipage}}

\begin{Shaded}
\begin{Highlighting}[]
\NormalTok{A.C }\OtherTok{=} \FunctionTok{interaction}\NormalTok{(A,C)}
\FunctionTok{interaction.plot}\NormalTok{(B,A.C, fits.r)}
\end{Highlighting}
\end{Shaded}

\includegraphics{hw6_files/figure-latex/unnamed-chunk-14-1.pdf}

\hypertarget{maximum-setting}{%
\paragraph{Maximum setting}\label{maximum-setting}}

In this characterization experiment, we find that the best combination
of factor levels is given by a low cutting speed (\(A^-\)), a high tool
geometry (\(B^+\)), and a high cutting angle (\(C^+\)).

\hypertarget{explanation}{%
\paragraph{Explanation}\label{explanation}}

First, we observe that \(A^-\) and \(C^+\) has the best interaction
effect (its corresponding line is above all of the others). Second, we
see that \(B\) has a positive effect, i.e., as we go from \(B^-\) to
\(B^+\), the response increases, and thus \(B^+\) is better. Taking
these two facts together, we see that the best combination in this
experiment is \((A^-,B^+,C^+)\).

We may also just look at the highest point on the graph and see that
\((A^-,B^+,C^+)\) has the highest mean response. However, the
interaction plot reveals more, particularly if in follow-up experiments
where we use what we learned in this experiment to narrow down our
search for an optimal combination (where optimality may be with respect
to, say, mean response or minimal variability in the response).

\hypertarget{appendix}{%
\section{Appendix}\label{appendix}}

\hypertarget{problem-1-additional-analysis-of-the-additive-model}{%
\subsection{Problem 1: additional analysis of the additive
model}\label{problem-1-additional-analysis-of-the-additive-model}}

The model parameter estimates are given by:

\begin{Shaded}
\begin{Highlighting}[]
\FunctionTok{dummy.coef}\NormalTok{(mod.R)}
\end{Highlighting}
\end{Shaded}

\begin{verbatim}
## Full coefficients are 
##                                             
## (Intercept):         3.125                  
## A:                      10         12     14
##                     -3.625     -0.625  4.250
## B:                      25         30       
##                     -1.375      1.375       
## C:                     200        250       
##                 -0.9583333  0.9583333
\end{verbatim}

Thus, the estimated mean response \(\hat{Y}_{i j k}^{(R)}\) is given by
\begin{align*}
  \hat{Y}_{i j k}^{(R)} = 3.125
    - &I(i=1) 3.625 - I(i=2) 6.25 + I(i=3) 4.250\\
    - &I(j=1) 1.375 + I(j=2) 1.375\\
    - &I(k=1) 9.58\bar{3} + I(k=2) 9.58\bar{3}.
\end{align*}

In R code, this may be brute-force implemented as:

\begin{Shaded}
\begin{Highlighting}[]
\NormalTok{yhat.R }\OtherTok{=} \ControlFlowTok{function}\NormalTok{(i,j,k)}
\NormalTok{\{}
\NormalTok{  y }\OtherTok{=} \FloatTok{3.125}
  \ControlFlowTok{if}\NormalTok{ (i }\SpecialCharTok{==} \DecValTok{1}\NormalTok{) y }\OtherTok{=}\NormalTok{ y }\SpecialCharTok{{-}} \FloatTok{3.625}\NormalTok{; }\ControlFlowTok{if}\NormalTok{ (i }\SpecialCharTok{==} \DecValTok{2}\NormalTok{) y }\OtherTok{=}\NormalTok{ y }\SpecialCharTok{{-}}\NormalTok{ .}\DecValTok{625}\NormalTok{; }\ControlFlowTok{if}\NormalTok{ (i }\SpecialCharTok{==} \DecValTok{3}\NormalTok{) y }\OtherTok{=}\NormalTok{ y }\SpecialCharTok{+} \FloatTok{4.25}
  \ControlFlowTok{if}\NormalTok{ (j }\SpecialCharTok{==} \DecValTok{1}\NormalTok{) y }\OtherTok{=}\NormalTok{ y }\SpecialCharTok{{-}} \FloatTok{1.375}\NormalTok{; }\ControlFlowTok{if}\NormalTok{ (j }\SpecialCharTok{==} \DecValTok{2}\NormalTok{) y }\OtherTok{=}\NormalTok{ y }\SpecialCharTok{+} \FloatTok{1.375}
  \ControlFlowTok{if}\NormalTok{ (k }\SpecialCharTok{==} \DecValTok{1}\NormalTok{) y }\OtherTok{=}\NormalTok{ y }\SpecialCharTok{{-}}\NormalTok{ .}\DecValTok{95833333}\NormalTok{; }\ControlFlowTok{if}\NormalTok{ (k }\SpecialCharTok{==} \DecValTok{2}\NormalTok{) y }\OtherTok{=}\NormalTok{ y }\SpecialCharTok{+}\NormalTok{ .}\DecValTok{95833333}
\NormalTok{  y}
\NormalTok{\}}
\end{Highlighting}
\end{Shaded}

We print out all of the estimates of the mean response with:

\begin{Shaded}
\begin{Highlighting}[]
\ControlFlowTok{for}\NormalTok{ (i }\ControlFlowTok{in} \DecValTok{1}\SpecialCharTok{:}\DecValTok{3}\NormalTok{)}
  \ControlFlowTok{for}\NormalTok{ (j }\ControlFlowTok{in} \DecValTok{1}\SpecialCharTok{:}\DecValTok{2}\NormalTok{)}
    \ControlFlowTok{for}\NormalTok{ (k }\ControlFlowTok{in} \DecValTok{1}\SpecialCharTok{:}\DecValTok{2}\NormalTok{)}
      \FunctionTok{cat}\NormalTok{(}\StringTok{"Y("}\NormalTok{,}\FunctionTok{c}\NormalTok{(i,j,k),}\StringTok{") = "}\NormalTok{,}\FunctionTok{yhat.R}\NormalTok{(i,j,k),}\StringTok{"}\SpecialCharTok{\textbackslash{}n}\StringTok{"}\NormalTok{)}
\end{Highlighting}
\end{Shaded}

\begin{verbatim}
## Y( 1 1 1 ) =  -2.833333 
## Y( 1 1 2 ) =  -0.9166667 
## Y( 1 2 1 ) =  -0.08333333 
## Y( 1 2 2 ) =  1.833333 
## Y( 2 1 1 ) =  0.1666667 
## Y( 2 1 2 ) =  2.083333 
## Y( 2 2 1 ) =  2.916667 
## Y( 2 2 2 ) =  4.833333 
## Y( 3 1 1 ) =  5.041667 
## Y( 3 1 2 ) =  6.958333 
## Y( 3 2 1 ) =  7.791667 
## Y( 3 2 2 ) =  9.708333
\end{verbatim}

\hypertarget{problem-1-plots-of-interactions-that-were-deemed-insignificant}{%
\subsection{Problem 1: plots of interactions that were deemed
insignificant}\label{problem-1-plots-of-interactions-that-were-deemed-insignificant}}

The data showed no statistically significant interaction effects between
factors \(A\) and \(B\) and between factors \(B\) and \(C\). Thus, their
interaction plots are, as expected, parallel:

\begin{Shaded}
\begin{Highlighting}[]
\FunctionTok{par}\NormalTok{(}\AttributeTok{mfrow=}\FunctionTok{c}\NormalTok{(}\DecValTok{1}\NormalTok{,}\DecValTok{2}\NormalTok{))}
\FunctionTok{interaction.plot}\NormalTok{(A,B,fits.r,}\AttributeTok{ylab=}\StringTok{"y"}\NormalTok{)}
\FunctionTok{interaction.plot}\NormalTok{(B,C,fits.r,}\AttributeTok{ylab=}\StringTok{"y"}\NormalTok{)}
\end{Highlighting}
\end{Shaded}

\includegraphics{hw6_files/figure-latex/unnamed-chunk-18-1.pdf}

\end{document}
